% paper_16_periodicity.tex — GeoVac Paper 16
% Chemical Periodicity as S_N Representation Theory on S^{3N-1}
% Status: Draft, March 18, 2026
\documentclass[aps,prl,twocolumn,superscriptaddress,longbibliography]{revtex4-2}

\usepackage{amsmath,amssymb,graphicx,xcolor,bm,braket,dcolumn,booktabs}
\usepackage{hyperref}

\newcommand{\Nelectron}{N}
\newcommand{\mufree}{\mu_{\mathrm{free}}}
\newcommand{\neff}{n_{\mathrm{eff}}}
\newcommand{\leff}{l_{\mathrm{eff}}}
\newcommand{\Zeff}{Z_{\mathrm{eff}}}
\newcommand{\Ha}{\,\mathrm{Ha}}
\newcommand{\eV}{\,\mathrm{eV}}

\begin{document}

\title{The Periodic Table Recast as $S_N$ Representation Theory\\
on $S^{3N-1}$}
\thanks{Paper 16 in the Geometric Vacuum series}

\author{J.~Loutey}
\affiliation{Independent Researcher, Kent, Washington}

\date{March 18, 2026}

%% ========================================================================
%% ABSTRACT
%% ========================================================================

\begin{abstract}
The non-relativistic $N$-electron Schr\"odinger equation, after
separation of the hyperradius, reduces to angular dynamics on the
hypersphere $S^{3N-1}$.  The angular kinetic energy is the
$\mathrm{SO}(3N)$ Casimir operator with eigenvalue
$\mufree = \nu(\nu + 3N - 2)/2$.  Fermionic antisymmetry restricts
the spatial wavefunction to a specific $S_N$ irreducible
representation---the conjugate of the spin Young diagram
$[n_\uparrow, n_\downarrow]$, a diagram with two columns of heights
$n_\uparrow$ and $n_\downarrow$.  The minimum angular quantum number
compatible with it is the Pauli filling count
$\nu_{\min} = E(n_\uparrow) + E(n_\downarrow)$, where $E(m)$ sums the
$m$ lowest one-particle polynomial degrees and degree $d$ holds
$(d+1)(d+2)/2$ polynomials.  We prove this for all $N$ and $S$ from
Schur--Weyl duality (it extends the known filled-shell rule of the
hyperspherical Fermi-gas literature to open shells and arbitrary spin)
and corroborate it by an exact character count for all spins through
$N = 12$ and by brute force for $N = 3, 4$.  It depends on the spin,
equals $N-2$ exactly when both columns lie in the first two oscillator
shells ($1 \le n_\downarrow \le n_\uparrow \le 4$; at lowest spin, for
$2 \le N \le 8$), and grows as $N^{4/3}$; the ground-state orbital
angular momentum and parity can raise it further (lithium's ${}^2S$
needs $\nu = 2$, not~1).  The $S_N$ irrep structure yields a
classification of atoms into five types---A (trivial), B~(noble gas),
C~(alkali), D~(alkaline earth), E~(open shell)---that recasts the periodic
table's group structure in representation-theoretic language.  The periodic law
corresponds to the repetition of the $S_N$ irrep sequence
$\mathrm{C} \to \mathrm{D} \to \mathrm{E} \to \mathrm{B}$ within each
period.  This classification maps onto, rather than predicts, the
known periodic table structure; its value lies in providing a
unified group-theoretic framework for a pattern that is usually
described empirically.  (An earlier version of this paper claimed a
universal $\nu = N-2$, $\mufree = 2(N-2)(N-1)$, and a limit of
4~units per electron pair; all three are withdrawn,
Sec.~\ref{sec:nu_gs}.)  We show that the topological
structure ($\nu$, $\mufree$) is smooth through $Z = 1/\alpha \approx
137$ and that the Dirac instability is a \emph{metric} singularity on
$S^{3N-1}$, not a topological one.
\end{abstract}

\maketitle

%% ========================================================================
%% I. INTRODUCTION
%% ========================================================================

\section{Introduction}
\label{sec:intro}

The periodic table of elements, established empirically by Mendeleev
in 1869~\cite{Mendeleev1869}, received its first theoretical
justification from quantum mechanics through the shell model and the
aufbau principle.  The $2n^2$ orbital counting rule, the grouping of
elements by valence shell, and the progression of ionization energies
across a period are all consequences of the hydrogen-like
Schr\"odinger equation and its multi-electron extensions.

Despite this success, the periodic law itself---the recurrence of
chemical properties with increasing atomic number---has not been
derived from a single mathematical principle.  Shell-filling rules are
empirical (the Madelung rule has exceptions), Hund's rules require
separate justification, and the distinction between noble gases, alkali
metals, and transition metals is described rather than predicted.

In this paper, we show that the periodic table's group structure
can be recast in the language of $S_N$ representation theory acting
on the spatial part of the $N$-electron wavefunction on the
hypersphere $S^{3N-1}$.  The key results are:

\begin{enumerate}
  \item The minimum angular quantum number compatible with the spin
    is the Pauli filling count $\nu_{\min} = E(n_\uparrow) +
    E(n_\downarrow)$, which depends on $N$ and $S$ but not on $Z$
    (Sec.~\ref{sec:framework}).
  \item The Pauli centrifugal cost $\mufree = \nu_{\min}(\nu_{\min} +
    3N-2)/2$ grows as $N^{8/3}$, faster than the number of electron
    pairs (Sec.~\ref{sec:mu_formula}).
  \item The shape of the $S_N$ irrep classifies atoms into chemical
    types that recast the group structure of the periodic table in
    representation-theoretic language
    (Sec.~\ref{sec:classification}).
  \item The periodic law corresponds to the repetition of $S_N$ irrep
    sequences (Sec.~\ref{sec:periodic_law}).
  \item The Dirac instability at $Z\alpha \to 1$ is a metric
    singularity, not a topological one; the representation theory
    is smooth through $Z = 137$
    (Sec.~\ref{sec:topology_vs_metric}).
\end{enumerate}

%% ========================================================================
%% I-B. WHAT IS NEW VS. WHAT IS KNOWN
%% ========================================================================

\section{What is New vs.\ What is Known}
\label{sec:new_vs_known}

Several ingredients of the present analysis are standard; it is
important to distinguish them from the new contributions.

\subsection{Known results}

\begin{itemize}
  \item \textbf{Hyperspherical formulation.}  The separation of the
    $N$-electron Schr\"odinger equation into hyperradial and angular
    parts on $S^{3N-1}$, and the identification of the angular kinetic
    energy with the $\mathrm{SO}(3N)$ Casimir, are
    established~\cite{Avery2000, Fock1954}.

  \item \textbf{$S_N$ representation theory for fermions.}  The use of
    Young diagrams and Schur--Weyl duality to classify spatial--spin
    factorizations of antisymmetric wavefunctions is textbook
    material~\cite{Weyl1946, FultonHarris1991}.

  \item \textbf{Minimal hyperangular momentum of filled oscillator
    shells.}  For $N$ non-interacting fermions in an isotropic
    oscillator, the ground state is a hyperspherical harmonic whose
    degree equals the total number of oscillator quanta; for filled
    shells this is stated and proved as Theorem~1 of
    Rittenhouse, Cavagnero, von Stecher and Greene~\cite{Rittenhouse2006},
    together with the large-$N$ law $(3N)^{4/3}/4$.  The related nuclear
    method of $K$ harmonics compares such hyperspherical states with the
    translation-invariant shell model~\cite{SmirnovShitikova1977} (there
    in $3(A-1)$ Jacobi coordinates); we have read only that review's
    abstract and do not attribute the rule to it.

  \item \textbf{The Casimir eigenvalue formula.}  The expression
    $\mufree = \nu(\nu + 3N - 2)/2$ follows directly from the
    standard $\mathrm{SO}(3N)$ representation theory.

  \item \textbf{The periodic table itself.}  The classification of
    elements into noble gases, alkali metals, alkaline earths, and
    transition metals is an empirical fact that guided the
    construction presented here.
\end{itemize}

\subsection{New contributions}

\begin{enumerate}
  \item \textbf{The minimum angular quantum number.}  The minimum
    $\nu$ compatible with the spin is the Pauli filling count of
    Eq.~(\ref{eq:nu_min}) (Theorem~2), together with the effect of the
    ground-state $L$ and parity (Table~\ref{tab:nu_verified}).  The
    filling count is the oscillator-quanta count of $N$ non-interacting
    fermions in a three-dimensional isotropic harmonic oscillator; for
    filled shells the rule is known~\cite{Rittenhouse2006}.  The
    contribution here is the spin-resolved, open-shell statement as an
    $S_N$-irrep result with a two-line Schur--Weyl proof, and its
    tabulation for atomic ground states on $S^{3N-1}$.  (An earlier
    version claimed a universal $\nu = N-2$ as new; that claim is
    withdrawn.)

  \item \textbf{The five-type classification (A/B/C/D/E).}  The
    observation that the $S_N$ irrep shape naturally partitions atoms
    into types that map onto the periodic table's group structure
    provides a group-theoretic reformulation of the empirical pattern.

  \item \textbf{The Dirac instability as metric vs.\ topological
    singularity.}  The analysis showing that $\nu$, $\mufree$, and
    the $S_N$ irrep are smooth through $Z = 137$ while the metric
    diverges clarifies the nature of the upper limit on the periodic
    table.
\end{enumerate}

%% ========================================================================
%% III. THE HYPERSPHERICAL FRAMEWORK
%% ========================================================================

\section{The Hyperspherical Framework}
\label{sec:framework}

\subsection{Configuration space and angular dynamics}
\label{sec:config}

Consider $N$ electrons in a central potential $V(r_1, \ldots, r_N)$.
After extracting the center-of-mass motion, the configuration space is
$\mathbb{R}^{3N}$.  Introducing hyperspherical coordinates
$(R, \hat{\Omega})$ where $R = \sqrt{\sum_i r_i^2}$ is the
hyperradius and $\hat{\Omega} \in S^{3N-1}$ collects $3N-1$ angular
variables, the kinetic energy separates as~\cite{Avery2000,
Fock1954}:
\begin{equation}
  T = -\frac{1}{2}\left(\frac{\partial^2}{\partial R^2}
    + \frac{3N-1}{R}\frac{\partial}{\partial R}
    - \frac{\Lambda^2}{R^2}\right),
  \label{eq:kinetic}
\end{equation}
where $\Lambda^2$ is the grand angular momentum operator---the
Laplace--Beltrami operator on $S^{3N-1}$.  Its eigenvalues are the
$\mathrm{SO}(3N)$ Casimir invariants:
\begin{equation}
  \Lambda^2 Y_\nu = \nu(\nu + 3N - 2)\, Y_\nu \equiv 2\mufree\, Y_\nu,
  \label{eq:casimir}
\end{equation}
with $\nu = 0, 1, 2, \ldots$ and hyperspherical harmonics $Y_\nu$.
The effective angular momentum quantum numbers are
\begin{equation}
  \leff = \frac{3N - 3}{2}, \quad \neff = \frac{3N - 1}{2}.
  \label{eq:eff_qn}
\end{equation}

The centrifugal barrier in the hyperradial equation is
$\leff(\leff + 1)/R^2 = (3N-3)(3N-1)/(4R^2)$, showing that even in
the lowest angular channel, $N$-electron atoms face a centrifugal cost
that grows as $\sim 9N^2/4$.

\subsection{Fermionic antisymmetry and $S_N$ irreps}
\label{sec:sn_irreps}

The total wavefunction must be antisymmetric under $S_N$ (permutation
of electrons).  Since the total state factorizes as $|\Psi\rangle =
|\text{spatial}\rangle \otimes |\text{spin}\rangle$, the spatial and
spin parts must carry conjugate $S_N$ irreducible representations.

For $N$ spin-1/2 fermions with total spin $S$, the spin irrep is a
Young diagram with at most two rows: $[n_\uparrow, n_\downarrow]$
where $n_\uparrow = N/2 + S$ and $n_\downarrow = N/2 - S$.  By
Schur--Weyl duality~\cite{Weyl1946, FultonHarris1991}, the spatial
irrep is the \emph{conjugate} (transposed) Young diagram.

\noindent\textbf{Lemma 1} (Two-column spatial diagram).
\label{thm:lambda1}
\textit{For an $N$-electron state of total spin $S$, the spatial Young
diagram has at most two columns, of heights $n_\uparrow = N/2 + S$ and
$n_\downarrow = N/2 - S$.  Equivalently, it has $n_\downarrow$ rows of
length~2 and $2S$ rows of length~1.}

\medskip\noindent\textit{Proof.}
The spin diagram $[n_\uparrow, n_\downarrow]$ has at most two rows
(spin-$\tfrac12$ Schur--Weyl duality); transposition exchanges rows and
columns. \hfill$\square$

\noindent For $S < N/2$ the first \emph{row} has length~2; the diagram's
first \emph{column} has height $n_\uparrow$, which does depend on $S$.
For $S = N/2$ the spatial diagram is the single column $[1^N]$: the
fully \emph{antisymmetric} irrep.
% [retracted 2026-09-28: p16-nu-universal-n-minus-2] The former
% Theorem 1 ("Universal lambda_1") stated that for S < N/2 the spatial
% diagram has "first-column length lambda_1 = 2", and treated S = N/2 as
% "lambda_1 = N, nu = 0 ... completely symmetric in space ... the
% Laughlin-like state with no Pauli nodes". The S = N/2 spatial state is
% fully antisymmetric and has nu_min = E(N) > 0.

\subsection{The minimum angular quantum number}
\label{sec:nu_gs}

The minimum $\nu$ compatible with a given $S_N$ irrep is the lowest
degree $K$ of a harmonic homogeneous polynomial on $\mathbb{R}^{3N}$
(a hyperspherical harmonic of degree~$K$ on $S^{3N-1}$) whose $S_N$
content contains that irrep.  Order the one-particle monomials on
$\mathbb{R}^3$ by degree; degree $d$ holds $g_d = (d+1)(d+2)/2 = 1, 3,
6, 10, \ldots$ of them, and let $E(m)$ be the sum of the degrees of the
$m$ lowest, i.e.
\begin{equation}
  E(m) = 3\binom{d+2}{4} + d\,\bigl(m - \tbinom{d+2}{3}\bigr),
  \qquad \tbinom{d+2}{3} \le m < \tbinom{d+3}{3}.
  \label{eq:E_of_m}
\end{equation}

\noindent\textbf{Theorem 2} (Pauli filling count).
\textit{For every $N \ge 1$ and every total spin $S$, the lowest degree
of a harmonic polynomial on $\mathbb{R}^{3N}$ carrying the spatial
irrep of Lemma~\ref{thm:lambda1} is}
\begin{equation}
  \nu_{\min}(N, S) = E(n_\uparrow) + E(n_\downarrow).
  \label{eq:nu_min}
\end{equation}

\medskip\noindent\textit{Proof.}
Let $W = \mathbb{C}[x,y,z] = \bigoplus_d W_d$ be the graded
one-particle polynomial space, $\dim W_d = g_d$.  The polynomials on
$\mathbb{R}^{3N}$ are $W^{\otimes N}$, with $S_N$ permuting the
factors.  By Schur--Weyl duality (applied degree by degree to the
finite-dimensional truncations)~\cite{FultonHarris1991}, the
multiplicity space of the $S_N$ irrep $\lambda$ in $W^{\otimes N}$ is
the Schur module $\mathbb{S}_\lambda W$, whose graded dimension is the
Schur function $s_\lambda$ evaluated on the degree grading, i.e.\ a sum
over semistandard tableaux $T$ of shape $\lambda$ filled with basis
monomials of $W$ of $t^{\deg T}$, with $\deg T$ the total degree of the
entries.  Every coefficient is a nonnegative count, so the lowest
degree at which $\lambda$ occurs is $\min_T \deg T$, with no
cancellation.  Entries strictly increase down each column, so a column
of height $h$ holds $h$ distinct monomials and has degree at least
$E(h)$; filling every column with the lowest monomials in order
(rows then weakly increase) attains this bound.  For the two-column
spatial diagram of Lemma~\ref{thm:lambda1}, of heights $n_\uparrow$
and $n_\downarrow$, the lowest degree is therefore
$E(n_\uparrow) + E(n_\downarrow)$.  Finally, polynomials of degree $K$
decompose as harmonics of degree $K$ plus $\bigl(\sum_i r_i^2\bigr)$
times polynomials of degree $K-2$, and $\sum_i r_i^2$ is
$S_N$-invariant; since $\lambda$ is absent in degree $K-2$ at the
lowest degree $K$, it occurs there among the harmonics.
\hfill$\square$

\noindent Equivalently, $\nu_{\min}$ is the minimum number of oscillator
quanta of $N$ non-interacting spin-$\tfrac12$ fermions of spin $S$ in a
three-dimensional isotropic oscillator, filled separately for the two
spin columns; for filled shells this is Theorem~1 of
Ref.~\cite{Rittenhouse2006}.  \emph{Evidence tier: symbolic proof
(above), corroborated computationally}: (A)~the exact character
generating function $(1-t^2)\,\frac{1}{N!}\sum_{\sigma}
\chi_{\tilde\lambda}(\sigma) \prod_{c}(1-t^{|c|})^{-3}$ (product over
the cycles of $\sigma$; characters by Murnaghan--Nakayama) reproduces
Eq.~(\ref{eq:nu_min}) for every spin through $N = 12$; (B)~the rank of
the Young symmetrizer on explicit monomials, an independent route with
no characters, agrees for $N = 3, 4$.  Backing:
\texttt{tests/test\_paper16\_nu\_min.py}.

Three consequences.  (i)~$\nu_{\min}$ depends on $N$ and $S$ but not
on $Z$.  (ii)~Since $E(m) = m-1$ exactly for $1 \le m \le 4$ (the first
two shells) and $E(m) > m-1$ for $m = 0$ and $m \ge 5$, $\nu_{\min} =
N-2$ holds exactly when $1 \le n_\downarrow \le n_\uparrow \le 4$.  At
the lowest spin that is $2 \le N \le 8$; it also covers higher spins
in that range (e.g.\ carbon's ${}^3P$, $N = 6$, $S = 1$).  From $N = 9$
at lowest spin each electron in the third shell costs two units, in
the fourth three, and $\nu_{\min} \simeq (3N)^{4/3}/4$ at large
$N$~\cite{Rittenhouse2006}.  (iii)~Raising the spin at fixed $N$ never
lowers $\nu_{\min}$ and generally raises it (nitrogen: 5, 6, 7, 9 for
$S = \tfrac12, \tfrac32, \tfrac52, \tfrac72$).  \textit{Proof of
(iii):} $E(m+1) - E(m)$ is the degree of the $(m+1)$-th monomial, which
is nondecreasing in $m$ ($E$ is convex).  Raising $S$ by one moves one
box from the $n_\downarrow$ column to the $n_\uparrow$ column, changing
$\nu_{\min}$ by $[E(n_\uparrow + 1) - E(n_\uparrow)] -
[E(n_\downarrow) - E(n_\downarrow - 1)] \ge 0$, because $n_\uparrow + 1
> n_\downarrow$.

The $S_N$ count ignores the ground state's orbital angular momentum
$L$ and parity.  Imposing them (replace $(1-t^c)^{-3}$ by
$\det(1 - t^c g^c)^{-1}$ and average over $g \in \mathrm{O}(3)$) can
raise $\nu_{\min}$ further: the degree-1 harmonics carrying lithium's
$[2,1]$ irrep are the coordinates themselves, $L = 1$ and odd parity
(${}^2P$), so the ${}^2S$ ground state needs $\nu = 2$.
Table~\ref{tab:nu_verified} lists both counts for $N = 2$--$12$.

\begin{table}[h]
\centering
\caption{Minimum angular quantum number for the Hund's-rule ground
terms, $N = 2$--$12$.  $\nu_{\min}^{S_N}$: spin (spatial irrep) only,
Eq.~(\ref{eq:nu_min}); $\nu_{\min}^{L,\pi}$: additionally the ground
term's $L$ and parity (exact $S_N \times \mathrm{O}(3)$ character
count).  The last column is the withdrawn rule $N-2$.}
\label{tab:nu_verified}
\begin{tabular}{llccc}
\toprule
Atom & Term & $\nu_{\min}^{S_N}$ & $\nu_{\min}^{L,\pi}$ & $N-2$ \\
\midrule
He & ${}^1S^e$ &  0 &  0 &  0 \\
Li & ${}^2S^e$ &  1 &  2 &  1 \\
Be & ${}^1S^e$ &  2 &  2 &  2 \\
B  & ${}^2P^o$ &  3 &  3 &  3 \\
C  & ${}^3P^e$ &  4 &  4 &  4 \\
N  & ${}^4S^o$ &  6 &  7 &  5 \\
O  & ${}^3P^e$ &  7 &  8 &  6 \\
F  & ${}^2P^o$ &  8 &  9 &  7 \\
Ne & ${}^1S^e$ & 10 & 10 &  8 \\
Na & ${}^2S^e$ & 12 & 12 &  9 \\
Mg & ${}^1S^e$ & 14 & 14 & 10 \\
\bottomrule
\end{tabular}
\end{table}

% [retracted 2026-09-28: p16-nu-universal-n-minus-2] Former text:
% "For the spatial irrep with first-column length lambda_1, the minimum
% angular excitation is [Avery2000]: nu_min = N - lambda_1. By Theorem 1,
% all atoms with S < N/2 have \boxed{nu = N - 2} as their ground-state
% angular quantum number. This is a universal result: it depends only on
% the electron count N, not on the nuclear charge Z, the spin S (provided
% S < N/2), or any details of the electron configuration."
% The Avery attribution is also withdrawn: the rule N - lambda_1 is false
% for N >= 9 (and spin-dependent from N = 7), so it cannot be sourced.

%% ========================================================================
%% III. THE mu_free FORMULA
%% ========================================================================

\section{The $\mufree$ Formula}
\label{sec:mu_formula}

Substituting $\nu_{\min}$ of Eq.~(\ref{eq:nu_min}) into
Eq.~(\ref{eq:casimir}):
\begin{equation}
  \mufree = \frac{\nu_{\min}(\nu_{\min} + 3N - 2)}{2}.
  \label{eq:mu_formula}
\end{equation}

This is the \emph{Pauli centrifugal cost}: the minimum angular kinetic
energy imposed by fermionic antisymmetry (at spin $S$) on $S^{3N-1}$.
It reduces to $2(N-2)(N-1)$ only where $\nu_{\min} = N-2$, i.e.\ when
$1 \le n_\downarrow \le n_\uparrow \le 4$ (at lowest spin,
$2 \le N \le 8$).
% [retracted 2026-09-28: p16-nu-universal-n-minus-2] Former boxed result
% mu_free = 2(N-2)(N-1) for all S < N/2, with limits mu_free/N^2 -> 2,
% mu_free/binom(N,2) -> 4 ("each antisymmetric electron pair contributes
% exactly 4 units of angular momentum", eq:per_pair), and
% nu/(3N-2) -> 1/3 "monotonically from below".

\subsection{Asymptotic behavior}
\label{sec:asymptotic}

Several limits of Eq.~(\ref{eq:mu_formula}) are physically
informative:

\begin{enumerate}
  \item \textbf{Growth.}  With $\nu_{\min} \simeq (3N)^{4/3}/4$,
    $\mufree \simeq \nu_{\min}^2/2$ grows as $N^{8/3}$: faster than
    the $N^2$ of the number of electron pairs.  The Pauli cost per
    pair, $\mufree/\binom{N}{2}$, therefore grows without bound (as
    $N^{2/3}$); it does not approach a constant.

  \item \textbf{Excitation fraction.}  The ratio $\nu_{\min}/(3N-2)$
    also grows as $N^{1/3}$ (0.50 at Ar, 1.30 at Og); the ground
    state does not sit at a fixed fraction of the sphere's natural
    angular scale.

  \item \textbf{Fractional ionization cost.}  Removing one electron
    changes the sphere dimension by $3/(3N-1) \to 0$.  For heavy
    atoms, ionization is a \emph{small perturbation} to the angular
    topology.
\end{enumerate}

Table~\ref{tab:mu_values} collects these quantities for atoms across
the periodic table.

\begin{table}[h]
\centering
\caption{Hyperspherical parameters for neutral atoms at the
Hund's-rule ground spin, $S_N$ (spin) constraint only:
$\nu_{\min}$ from Eq.~(\ref{eq:nu_min}), $\mufree$ from
Eq.~(\ref{eq:mu_formula}).  For $N \le 12$ the $\nu_{\min}$ values are
exact character counts (Table~\ref{tab:nu_verified}); for $N \ge 18$
they are Eq.~(\ref{eq:nu_min}) with $S = 0$.  Imposing the ground
term's $L$ and parity can raise $\nu_{\min}$ (Li: 2, $\mufree = 9$;
N: 7, $\mufree = 91$).}
\label{tab:mu_values}
\begin{tabular}{lrrrrrc}
\toprule
Atom & $N$ & $S^{3N-1}$ & $\nu_{\min}$ & $\mufree$ & $\mufree/N^2$ & Type \\
\midrule
H     &   1 & $S^2$    &   0 &       0 & 0.00 & A \\
He    &   2 & $S^5$    &   0 &       0 & 0.00 & B \\
Li    &   3 & $S^8$    &   1 &       4 & 0.44 & C \\
Be    &   4 & $S^{11}$ &   2 &      12 & 0.75 & D \\
N     &   7 & $S^{20}$ &   6 &      75 & 1.53 & E \\
Ne    &  10 & $S^{29}$ &  10 &     190 & 1.90 & B \\
Na    &  11 & $S^{32}$ &  12 &     258 & 2.13 & C \\
Mg    &  12 & $S^{35}$ &  14 &     336 & 2.33 & D \\
Ar    &  18 & $S^{53}$ &  26 &    1014 & 3.13 & B \\
Kr    &  36 & $S^{107}$&  78 &    7176 & 5.54 & B \\
Xe    &  54 & $S^{161}$& 146 &   22338 & 7.66 & B \\
Rn    &  86 & $S^{257}$& 290 &   79170 & 10.7 & B \\
Og    & 118 & $S^{353}$& 456 &  184224 & 13.2 & B \\
\bottomrule
\end{tabular}
\end{table}

%% ========================================================================
%% IV. STRUCTURE TYPE CLASSIFICATION
%% ========================================================================

\section{Structure Type Classification}
\label{sec:classification}

Beyond the numbers $\nu_{\min}$ and $\mufree$, which depend only on
$N$ and $S$, the \emph{shape} of the $S_N$ irrep varies with the
electron configuration.  This shape encodes the internal structure of the
atom---how the $N$~electrons partition into core and valence.

We define five structure types based on the spatial Young diagram:

\subsection{Type A: Trivial ($N = 1$)}

The hydrogen atom has a single electron.  The trivial irrep $[1]$ of
$S_1$ gives $\nu_{\min} = 0$, $\mufree = 0$.  There is no
Pauli constraint.  The wavefunction lives on $S^2$ (a single
electron's angular space).

\subsection{Type B: Noble gas (democratic structure)}

Noble gases have closed-shell ground states with $S = 0$.  The spin
diagram is $[N/2,\, N/2]$ and the spatial irrep is its conjugate:
$[2, 2, \ldots, 2]$ ($N/2$ rows of length~2).

This is the most \emph{democratic} irrep: all columns have the same
height.  No electron is distinguished.  The wavefunction has the full
point-group symmetry of the Young diagram, and ionization requires
breaking this symmetry---a costly operation reflected in the high
ionization energies of noble gases (He:~$24.6\eV$, Ne:~$21.6\eV$,
Ar:~$15.8\eV$).

\subsection{Type C: Alkali metal (hierarchical structure)}

Alkali metals have one valence electron outside a noble-gas core, with
$S = 1/2$.  The spin diagram is $[(N+1)/2, \,(N-1)/2]$ and the
spatial irrep is $[2, 2, \ldots, 2, 1]$ ($= [2^{(N-1)/2}, 1]$).

This is a \emph{hierarchical} irrep: one column (the valence electron)
is shorter than the rest.  The wavefunction has a natural
decomposition into a symmetric core and a weakly bound valence
electron.  Ionization removes the short column---a topologically clean
surgery on $S^{3N-1}$:
\begin{equation}
  S^{3N-1} \to S^{3(N-1)-1} \times \mathbb{R}.
  \label{eq:ionization_surgery}
\end{equation}
The low cost of this surgery explains the low ionization energies
of alkali metals (Li:~$5.4\eV$, Na:~$5.1\eV$, K:~$4.3\eV$).

\subsection{Type D: Alkaline earth (core + pair)}

Alkaline earth metals have two valence electrons outside a noble-gas
core, with $S = 0$.  The spatial irrep is $[2, 2, \ldots, 2]$ (same
shape as Type~B) but the \emph{occupancy pattern} distinguishes the
core from the two valence electrons.

The ionization energy is intermediate (Be:~$9.3\eV$, Mg:~$7.6\eV$,
Ca:~$6.1\eV$)---higher than alkali metals (the valence pair is harder
to break) but lower than noble gases (the valence pair is still
distinguishable from the core).

\subsection{Type E: Open shell (partially filled $p$ or $d$)}

Atoms with a partially filled $p$ shell (boron through fluorine) or
$d$ shell (the transition metals) carry an open valence shell with
$S > 0$ but $S < N/2$.  The spatial irrep is neither the democratic
$[2^{N/2}]$ of Type~B/D nor the single-defect $[2^{(N-1)/2},1]$ of
Type~C, but an intermediate shape with several columns of unequal
height set by the open-shell occupancy and Hund's rule.  This is the
broadest class and the one where intra-shell exchange coupling is
strongest; the composed-builder implementation refines it further into
a $p$-block dispatch (type~E) and a separate $d$-block dispatch
(type~F) because the angular machinery differs, but at the
representation-theory level both are open-shell instances of the same
intermediate-irrep family.  Ionization energies span the widest range
(B:~$8.3\eV$ to F:~$17.4\eV$), tracking the filling of the open shell.

Table~\ref{tab:types} summarizes the classification.

\begin{table}[h]
\centering
\caption{Structure type classification.  The spatial Young diagram is
the conjugate of the spin diagram.  IE$_1$ is the first ionization
energy (experimental).}
\label{tab:types}
\begin{tabular}{cllcr}
\toprule
Type & Spatial irrep & Chemistry & Examples & IE$_1$ range \\
\midrule
A & $[1]$       & Monatomic       & H            & $13.6\eV$ \\
B & $[2^{N/2}]$ & Noble gas       & He, Ne, Ar   & $15$--$25\eV$ \\
C & $[2^{(N-1)/2}, 1]$ & Alkali  & Li, Na, K    & $4$--$5\eV$ \\
D & $[2^{N/2}]$ & Alkaline earth  & Be, Mg, Ca   & $6$--$9\eV$ \\
E & Various     & Open shell      & B--F, d-block& $8$--$17\eV$ \\
\bottomrule
\end{tabular}
\end{table}

%% ========================================================================
%% V. THE PERIODIC LAW AS IRREP SEQUENCE
%% ========================================================================

\section{The Periodic Law as Irrep Sequence}
\label{sec:periodic_law}

The periodic table's group structure arises because the sequence of
$S_N$ structure types \emph{repeats} with each period.  Within a given
period (defined by the outermost principal quantum number $n$), the
element sequence is:
\begin{equation}
  \mathrm{C} \to \mathrm{D} \to \underbrace{\mathrm{E} \to \cdots \to
  \mathrm{E}}_{p\text{- and }d\text{-block}} \to \mathrm{B}.
  \label{eq:period_sequence}
\end{equation}
Each period begins with an alkali metal (Type~C: one electron in the
new shell) and ends with a noble gas (Type~B: closed shell).  The
intermediate elements (Type~E) have partially filled subshells with
varying spin and Young diagram shapes.

In group-theoretic language, the periodic law can be understood as
follows:
\begin{quote}
\emph{The chemical properties of elements are periodic functions of
atomic number because the sequence of $S_N$ spatial irreps repeats
whenever a new principal shell begins filling.}
\end{quote}
This is a reformulation, not a derivation from first principles: the
shell-filling order itself (which determines where the sequence
repeats) is an empirical input, not a consequence of $S_N$
representation theory alone.

The key distinction between Types~B and~D is worth emphasizing.
Both have the same Young diagram shape $[2^{N/2}]$, but they differ
in the \emph{occupancy pattern}: Type~B fills all available subshells,
while Type~D has two electrons in a new shell above a closed core.
This distinction is invisible at Level~1 (representation theory) and
becomes visible only at Level~2 (the perturbative energy, which
depends on $Z$, orbital angular momenta, and electron--electron
repulsion).

\subsection{Isoelectronic series}
\label{sec:isoelectronic}

Atoms with the same $N$ but different $Z$ share identical $S_N$ irreps
and structure types.  The 10-electron series
\begin{equation}
  \text{Ne},\ \text{Na}^+,\ \text{Mg}^{2+},\ \text{Al}^{3+},\ \ldots
\end{equation}
all have the Type~B spatial irrep $[2,2,2,2,2]$ with $\nu_{\min} = 10$
and $\mufree = 190$.  Their chemistry differs only through the nuclear
charge~$Z$, which enters at Level~2.  This cleanly separates
\emph{topology} (same for the series) from \emph{potential} (varies
with $Z$).

\subsection{Topology flow: ionization as surgery}
\label{sec:topology_flow}

Ionization changes the electron count $N \to N - 1$, transforming the
configuration space via Eq.~(\ref{eq:ionization_surgery}).  This is a
codimension-3 surgery that removes three dimensions (one electron's
$\mathbb{R}^3$) from the hypersphere.  The fractional dimension loss
is $3/(3N-1)$, which decreases monotonically:
\begin{equation}
  \frac{3}{3N - 1} \to 0 \quad (N \to \infty).
\end{equation}
For heavy atoms, ionization is a negligible topological perturbation.
This correlates with the empirical observation that ionization energies
decrease down a group: removing one electron from $S^{353}$ (Og) is
geometrically less disruptive than removing one from $S^5$ (He).

%% ========================================================================
%% VI. TWO-LEVEL SEPARATION
%% ========================================================================

\section{Topology versus Metric: Two-Level Structure}
\label{sec:topology_vs_metric}

The preceding analysis reveals a clean two-level separation:

\begin{itemize}
  \item \textbf{Level 1 (Topological).}  The quantities
    $\nu_{\min}$ and $\mufree$ in their $S_N$-only form
    (Eq.~(\ref{eq:nu_min})), the $S_N$ irrep, and the structure type
    depend only on the electron count $N$ and the total spin $S$.
    They are independent of $Z$ and of all details of the
    electron--electron interaction.  (Adding the ground term's $L$
    and parity, Table~\ref{tab:nu_verified}, is still
    $Z$-independent but no longer a function of $N$ and $S$ alone.)

  \item \textbf{Level 2 (Perturbative/Metric).}  The binding energy,
    ionization energy, electron configuration, and Hund's rules all
    depend on $Z$, the fine-structure constant $\alpha$, and the
    shape of the potential.  These quantities determine the
    \emph{chemistry} within a given topological class.
\end{itemize}

An important consequence: \emph{Hund's first rule is not a topological
effect.}  Level~1 does see the spin, but in the \emph{opposite}
direction: raising $S$ at fixed $N$ never lowers $\nu_{\min}$ and
generally raises it (Eq.~(\ref{eq:nu_min}); nitrogen's ${}^4S$ has
$\nu_{\min} = 6$ against~5 for a doublet).  The preference for
high-spin ground states therefore enters through the Coulomb exchange
integrals at Level~2, which must overcome a larger Level-1 angular
cost.
% [retracted 2026-09-28: p16-nu-universal-n-minus-2] Former text: "Since
% lambda_1 = 2 for all spatial irreps with S < N/2, the angular quantum
% number nu = N - 2 is independent of the spin value."

\subsection{The Dirac limit as metric singularity}
\label{sec:dirac_limit}

The Dirac equation for hydrogen-like atoms predicts that the 1s
binding energy diverges at $Z\alpha = 1$ ($Z \approx 137$ for a point
nucleus):
\begin{equation}
  E_{\mathrm{1s}}^{\mathrm{Dirac}}
    = mc^2 \left(\sqrt{1 - (Z\alpha)^2} - 1\right).
  \label{eq:dirac_1s}
\end{equation}
Does this ``Dirac dive'' have a topological signature on $S^{3N-1}$?

At $N = 137$ (lowest spin, $S = \tfrac12$): $\nu_{\min} = 570$ and
$\mufree = 279{,}015$ from Eqs.~(\ref{eq:nu_min})
and~(\ref{eq:mu_formula}).  These are integers given by a
piecewise-linear filling formula in $N$ whose only structure is the
oscillator shell closures ($N = 2, 8, 20, 40, 70, 112, 168, \ldots$ for
the singlet); nothing happens at $N = 137$.
% [retracted 2026-09-28: p16-nu-universal-n-minus-2] Former values at
% N = 137: nu = 135, mu_free = 36,720, nu/(3N-2) = 0.3301 "approaches
% 1/3 from below".

The Dirac instability is instead a \emph{metric} effect.  The
non-relativistic Fock projection yields the round metric on $S^3$.
Relativistic corrections warp this metric by a conformal factor
\begin{equation}
  \delta_{\mathrm{1s}} = \frac{(Z\alpha)^2}{1 - (Z\alpha)^2},
  \label{eq:metric_correction}
\end{equation}
which diverges at $Z\alpha \to 1$.  This ``pinch'' is a conical
singularity at the nuclear pole of $S^{3N-1}$---the geometric avatar
of the Dirac dive.  The topology ($S_N$ irrep, $\nu$, $\mufree$)
is oblivious; the metric (curvature, conformal factor) is singular.

Table~\ref{tab:metric} quantifies the transition from perturbative
to singular.

\begin{table}[h]
\centering
\caption{Relativistic metric correction $\delta_{1\mathrm{s}} =
(Z\alpha)^2/[1 - (Z\alpha)^2]$ for the innermost electron.  The
non-relativistic $\mufree$ formula becomes unreliable when
$\delta \gtrsim 1$.}
\label{tab:metric}
\begin{tabular}{rlrl}
\toprule
$Z$ & $(Z\alpha)^2$ & $\delta_{1\mathrm{s}}$ & Regime \\
\midrule
  1 & 0.00005 & 0.00005 & Non-relativistic \\
 29 & 0.045   & 0.047   & Mildly relativistic \\
 47 & 0.118   & 0.133   & Relativistic \\
 79 & 0.332   & 0.498   & Strongly relativistic \\
118 & 0.741   & 2.87    & Near-singular \\
137 & 0.999   & $\sim 10^3$ & Conical singularity \\
\bottomrule
\end{tabular}
\end{table}

\subsection{Three limits on the periodic table}

The periodic table terminates due to physics beyond the
non-relativistic $S_N$ representation theory:
\begin{enumerate}
  \item \textbf{Nuclear instability} ($Z \sim 126$): proton--proton
    Coulomb repulsion overwhelms the strong nuclear force.
  \item \textbf{Dirac instability, point nucleus} ($Z \sim 137$): the
    1s binding energy reaches $mc^2$.
  \item \textbf{Dirac instability, finite nucleus} ($Z \sim 170$):
    spontaneous $e^+e^-$ pair creation from the vacuum.
\end{enumerate}
The SO$(3N)$ topology predicts no limit: $\mufree$ of
Eq.~(\ref{eq:mu_formula}) is defined for all $N$.  The physical end of the periodic table is set by
nuclear physics ($Z \lesssim 130$), not by electronic topology.

%% ========================================================================
%% VII. DISCUSSION
%% ========================================================================

\section{Discussion}
\label{sec:discussion}

\subsection{Connection to Paper 0's orbital counting}

Paper~0~\cite{loutey_paper0} derived the $2n^2$ orbital degeneracy
from the geometric packing of nodes on the discrete $S^3$ lattice.
The present result extends this: the $S_N$ representation theory on
$S^{3N-1}$ yields not just the orbital counting but the
\emph{chemical classification} of atoms.  Where Paper~0 derived the $2n^2$ counting from geometric packing,
the present analysis provides a group-theoretic framework for
understanding why filling those shells produces the periodic pattern
$\mathrm{C} \to \mathrm{D} \to \mathrm{E} \to \mathrm{B}$.

\subsection{Transition metals and Hund's rule}

The transition metals (Type~E) have partially filled $d$-shells with
ground-state spins $S \geq 1$.  The ``anomalous'' configurations of
Cr~($[\mathrm{Ar}]\,3d^5\,4s^1$, $S = 3$) and
Cu~($[\mathrm{Ar}]\,3d^{10}\,4s^1$, $S = 1/2$) are often cited as
evidence for exchange stabilization.  In the present framework
Level~1 sees only $N$ and $S$: the ``expected'' $3d^4\,4s^2$ ($S = 2$)
has $\nu_{\min} = 42$ and the ``anomalous'' $3d^5\,4s^1$ ($S = 3$) has
$\nu_{\min} = 43$ (Eq.~(\ref{eq:nu_min})), so the topological cost
\emph{disfavors} the observed high-spin configuration.  The anomaly
is decided at Level~2, by exchange: the half-filled
$3d^5\,4s^1$ configuration has 141~exchange pairs versus 136 for
$3d^4\,4s^2$, a difference of~5.  The pair count itself depends only
on $(n_\uparrow, n_\downarrow)$; what is Level~2 is the exchange
energy each pair contributes, which tips the energetic balance
against the Level-1 cost.
% [retracted 2026-09-28: p16-nu-universal-n-minus-2] Former text: "both
% the expected and anomalous configurations of Cr have nu = 22 and
% mu_free = 1012 -- the topological Level 1 cannot distinguish them."

This confirms that Hund's first rule is a \emph{perturbative} effect
(exchange stabilization in the Coulomb interaction) rather than a
\emph{topological} one: the Level-1 angular cost on $S^{3N-1}$ rises,
not falls, with spin.
More broadly, the actual chemistry---which electron configuration wins
energetically---is determined by exchange, correlation, and
orbital-dependent Coulomb effects that operate below the resolution of
the topological classification.  The Level~1 framework identifies
which configurations are \emph{topologically equivalent}; it cannot
determine which one is the ground state.

\subsection{Physical chemistry correspondences}

The structure type classification connects naturally to familiar
chemical trends:
\begin{itemize}
  \item \textbf{Ionization energy} correlates with the topological
    cost of surgery on $S^{3N-1}$.  Type~B (noble gas) has the highest
    surgery cost; Type~C (alkali) has the lowest.
  \item \textbf{Atomic radius} correlates with the location of the
    surgery: Type~C atoms have a weakly bound outer electron (large
    radius), while Type~B atoms are compact.
  \item \textbf{Electronegativity} increases across a period
    ($\mathrm{C} \to \mathrm{B}$), mirroring the transition from
    hierarchical to democratic irreps.
  \item \textbf{The periodic law} can be understood as the
    repetition of the $S_N$ irrep sequence whenever a new principal
    shell begins to fill.
\end{itemize}

\subsection{Fractal structure: core and valence at every scale}

A striking feature of the classification is its \emph{scale
invariance}.  The core--valence decomposition at each ionization step
produces the same Types~B/C/D structure at every level:
\begin{align*}
  \text{Na} &\to \text{Na}^+ + e^-
    & (\text{C} &\to \text{B} + \text{A}), \\
  \text{Na}^+ &\to \text{Na}^{2+} + e^-
    & (\text{B} &\to \text{C} + \text{A}).
\end{align*}
Each ionization peels off one column from the Young diagram,
converting one type into another in a predictable sequence.  This
self-similar structure is a direct consequence of the recursive
nature of $S_N \supset S_{N-1}$ representations.

%% ========================================================================
%% VIII. CONCLUSION
%% ========================================================================

\section{Conclusion}
\label{sec:conclusion}

We have shown that the group structure of the periodic table---the
classification of elements into noble gases, alkali metals, alkaline
earths, and open-shell atoms---follows from the $S_N$ representation
theory of fermionic wavefunctions on the hypersphere $S^{3N-1}$.  The
key results are:

\begin{enumerate}
  \item The minimum angular quantum number compatible with spin $S$
    is the Pauli filling count $\nu_{\min} = E(n_\uparrow) +
    E(n_\downarrow)$ (Theorem~2, proved for all $N$ and $S$ by
    Schur--Weyl duality, extending the known filled-shell
    rule~\cite{Rittenhouse2006}), which equals $N-2$ exactly when
    $1 \le n_\downarrow \le n_\uparrow \le 4$ (at lowest spin,
    $2 \le N \le 8$); the ground
    term's $L$ and parity can raise it (Li ${}^2S$: 2).  The Pauli
    centrifugal cost $\mufree$ grows as $N^{8/3}$.

  \item The shape of the $S_N$ spatial irrep classifies atoms into
    types (A/B/C/D/E) that recast the chemical group structure in
    representation-theoretic language.

  \item The periodic law corresponds to the repetition of the irrep
    sequence $\mathrm{C} \to \mathrm{D} \to \mathrm{E} \to
    \mathrm{B}$ within each period.

  \item Hund's first rule and configuration anomalies (Cr, Cu) are
    Level~2 perturbative effects; the Level~1 cost rises with spin,
    so exchange must overcome it.

  \item The Dirac instability at $Z\alpha \to 1$ is a metric
    singularity on $S^{3N-1}$, not a topological one; the
    representation theory is smooth through all $N$.
\end{enumerate}

The approach separates the periodic table into a universal
\emph{topological skeleton} (Level~1: $\nu$, $\mufree$, structure
type) and a chemistry-specific \emph{metric flesh} (Level~2: $Z$,
$\alpha$, Coulomb integrals).  This separation provides a
group-theoretic framework for the periodic table's structure
(topology) while leaving the element-specific chemistry to
perturbative effects (metric).  The framework recasts, rather than
derives from first principles, the empirical periodic pattern---its
value lies in unifying the classification under a single
mathematical structure.

\begin{acknowledgments}
This work uses the hyperspherical formalism developed in the Geometric
Vacuum program (Papers~0--15).  The $S_N$ analysis was informed by the
classical references of Avery~\cite{Avery2000} and
Fock~\cite{Fock1935, Fock1954}.
\end{acknowledgments}

%% ========================================================================
%% REFERENCES
%% ========================================================================

\begin{thebibliography}{25}

\bibitem{Mendeleev1869}
D.~I.~Mendeleev,
\emph{Zh.\ Russ.\ Khim.\ Obshch.} \textbf{1}, 60 (1869).

\bibitem{Fock1935}
V.~Fock,
\emph{Z.\ Phys.} \textbf{98}, 145 (1935).

\bibitem{Fock1954}
V.~Fock,
\emph{Izv.\ Akad.\ Nauk SSSR, Ser.\ Fiz.} \textbf{18}, 161 (1954).

\bibitem{Avery2000}
J.~Avery,
\emph{Hyperspherical Harmonics and Generalized Sturmians}
(Kluwer, Dordrecht, 2000).

\bibitem{Rittenhouse2006}
S.~T.~Rittenhouse, M.~J.~Cavagnero, J.~von Stecher, and C.~H.~Greene,
``Hyperspherical description of the degenerate Fermi gas: $s$-wave
interactions,'' \emph{Phys.\ Rev.\ A} \textbf{74}, 053624 (2006),
arXiv:cond-mat/0607340.

\bibitem{SmirnovShitikova1977}
Yu.~F.~Smirnov and K.~V.~Shitikova, ``Method of $K$ harmonics and the
shell model,'' \emph{Sov.\ J.\ Part.\ Nucl.} \textbf{8}, 344 (1977).

\bibitem{Weyl1946}
H.~Weyl,
\emph{The Classical Groups: Their Invariants and Representations}
(Princeton University Press, Princeton, 1946).

\bibitem{FultonHarris1991}
W.~Fulton and J.~Harris,
\emph{Representation Theory: A First Course}
(Springer, New York, 1991).

\bibitem{loutey_paper0}
J.~Loutey,
``Angular Momentum as Information Geometry: Isotropic Packing and the Recurring Angular Cross-Section,''
GeoVac Paper~0 (2025).

\bibitem{loutey_paper1}
J.~Loutey,
``The Geometric Atom: the Hydrogen State Space as a Label Lattice,''
GeoVac Paper~1 (2025).

\bibitem{loutey_paper7}
J.~Loutey,
``The Dimensionless Vacuum: Recovering the Schr\"odinger Equation on a Scale-Invariant Label Lattice,''
GeoVac Paper~7 (2026).

\bibitem{loutey_paper13}
J.~Loutey,
``The Hyperspherical Lattice: Two-Electron Atoms as Coupled Channel Graphs,''
GeoVac Paper~13 (2026).

\end{thebibliography}

\end{document}
